\uuid{fW4O}
\titre{Variables continues}
\auteur{worms}
\organisation{exo7}
\licence{CC BY-NC-SA 4.0}
\datecreate{2017-08-28}
\niveau{L2}
\module{Probabilité et statistique}
\chapitre{Probabilité continue}
\sousChapitre{Densité de probabilité}
\difficulte{}
\theme{Variables continues}
\exo7qcmid{worms.worms.041}

\contenu{
\begin{question}
\qclassification{261.01}
On suppose que le temps d'attente (en minutes) d'un bus est une variable aléatoire de densité définie par
$$ f_X(x)= \frac{2}{15}\left(1-\frac{x}{15}\right)\mathbb{I}_{[0,15]}(x) $$
Si l'on suppose que ces temps d'attente pendant 10 jours sont indépendants et de même loi (celle définie précédemment), quelle est la probabilité qu'on ait, durant ces 10 jours, a attendre plus de 10 minutes au moins 3 fois ?
\begin{answers}
\bad{Environ $5\%$.}
\good{Environ $9\%$.}
\bad{Environ $21\%$.}
\bad{Environ $90\%$.}
\bad{On ne peut pas répondre à la question, il manque des éléments pour mener le calcul.}
\end{answers}
\begin{explanations}
On est dans un schéma binomial, on cherche la probabilité $\mathbb{P}(N\geq 3)$ où $N$ est de loi binomiale de paramètres $10$ et $p=\mathbb{P}(X>10)=1/9$ (après calcul). On trouve environ $\mathbb{P}(N\geq 3)=1-\sum_{k=0}^2 \mathbb{P}(N=k)\simeq 9\%$.
\end{explanations}
\end{question}
}
